\section{Introduction}

Coding agents have moved from completing lines to authoring change sets. Given
a repository and an issue report, an agent edits files, runs the test suite,
and iterates until the build is green~\cite{swe-bench}, with the interface it
works through setting how far it gets~\cite{swe-agent}. Code
written this way is entering ordinary projects. The quality question that
follows is procedural. On what basis is a change written by an agent admitted
into a system, and on what basis is it taken back out?

For code covered by tests the answer is familiar. A failing test names the
behavior that broke and narrows the search for the code responsible. Other
classes of generated code have no comparable signal. Text-to-SQL is one of
them. The task is to turn a natural-language question into a query over a
relational database, and the standard measure is execution accuracy, under
which a query counts as correct only when running it returns the same rows as a
reference query~\cite{bird}. A run over the development set of the BIRD
benchmark collapses 1,534 such verdicts into one percentage. Individual wrong
answers are often noise: the query reads the question correctly and disagrees
with the reference only on a convention---an extra projected column, a
\texttt{DISTINCT} that deduplicates a fan-out join, a \texttt{NULL} in the top
row of an ordered result. The defects worth writing code against are the ones
that recur across many questions, and the aggregate hides them. A score that
fails to move reports that something is wrong without reporting which code is
wrong or why.

Two things make that gap consequential. First, generated SQL fails quietly. A
query with a wrong join or a wrong aggregation returns a table rather than an
exception, so the defect travels to whoever reads the answer. Second, the
systems that emit such queries are no longer single model calls. The strongest are
multi-stage pipelines that wrap the model in schema linking, candidate
generation, selection, and
repair~\cite{chase-sql,xiyan-sql,metadata-extraction,agentar-scale-sql}. Every
stage is a piece of code whose only justification is the quality of the queries
that come out of it. Once an agent writes those stages, a project maintains code it
did not author, aimed at defects it never enumerated, and judged by a number
that cannot be traced back to any stage.

Existing pipelines report their resulting accuracy, while autonomous
research frameworks also evaluate proposed pipeline changes~\cite{ai-scientist-enterprise}.
Our focus is how per-question and per-step records support acceptance and
later rollback of agent-written code.

We used measured acceptance and rollback to develop a text-to-SQL pipeline. Humans designed an
instrumented harness and fixed a single rule of admission. Inside it a frontier coding agent
(Claude Opus, working in Claude Code) ran the development loop: run the
benchmark, mine the recorded run for a recurring failure signature, propose a
pipeline step aimed at that signature, and measure the net effect. Humans
reviewed proposals and could inspect the selected question group, reverted the steps whose gains
did not survive later runs, and wrote none of the pipeline code. We ask three
research questions.

\begin{itemize}
\item[RQ1] Can measured acceptance and rollback act as a practical quality
  gate for pipeline code that no human wrote, admitting a step only on its
  measured effect on the SQL the pipeline emits and removing it once that
  effect stops holding?
\item[RQ2] Which quality defects of generated SQL does trace-driven error
  mining surface, and what repair code does each justify?
\item[RQ3] Does the quality of an artifact accepted under that gate transfer to
  a different generating model, or does it belong to the model it was measured
  against?
\end{itemize}

The thesis of this paper is that execution records can support acceptance
and rollback of agent-written changes:
pipeline code written entirely by an agent is fit to keep when every step has
to earn its place by a measured effect on the executed artifact and loses that
place as soon as the effect stops holding.

The 17 steps that survived the gate raise execution accuracy on BIRD's
development set from 70.9\% for a single call to the same model to 76.2\%, and
reach 81.3\% on the private test set that no part of the development loop ever
touched. The paper contributes:

\begin{itemize}
\item a quality gate for agent-written code, stated as a procedure: propose
  from a failure signature localized in run records, measure over the closed
  set of questions the change can affect, re-measure over the full question
  set, and roll the change back when a later run argues against it;
\item an instrumented harness that keeps the gate cheap enough to run
  repeatedly: every generated query is executed and scored against a reference
  answer,
  each disagreement is classified row by row, every step reports its own
  decisions, and any two runs can be compared question by question;
\item the artifact the gate produced, a pipeline of 17 steps in five groups
  (table selection, schema building, generation and review, targeted repairs,
  normalizers), each bounding its own worst case and detecting its target with
  deterministic code wherever the schema or the query's syntax supplies the
  evidence, reported with its measured quality and with a stage-level
  decomposition of the gain read off ordinary runs rather than ablation
  reruns;
\item transfer measurements with other generating models, including a
  GigaCode improvement from 59.8\% to 66.5\% on BIRD development.
\end{itemize}
